\documentclass[12pt]{article}
\usepackage{amsmath,amssymb,amsthm,mathtools}
\usepackage{geometry}
\geometry{margin=1.2in}
\usepackage{hyperref}
\hypersetup{colorlinks=true,linkcolor=blue,citecolor=blue,urlcolor=blue}
\usepackage{booktabs}
\usepackage{array}
\usepackage{microtype}

% --- Theorem environments ---
\newtheorem{theorem}{Theorem}[section]
\newtheorem{corollary}[theorem]{Corollary}
\newtheorem{lemma}[theorem]{Lemma}
\newtheorem{proposition}[theorem]{Proposition}
\newtheorem{definition}[theorem]{Definition}
\newtheorem{remark}[theorem]{Remark}

% --- Shorthand ---
\newcommand{\F}{\mathbb{F}}
\newcommand{\Z}{\mathbb{Z}}
\newcommand{\N}{\mathbb{N}}
\newcommand{\R}{\mathbb{R}}
\newcommand{\C}{\mathbb{C}}
\newcommand{\O}{\mathbb{O}}
\newcommand{\OP}{\mathbb{O}P}
\newcommand{\qprime}{q}
\newcommand{\genus}{g}
\newcommand{\valency}{h}
\newcommand{\Roots}{\Phi}
\newcommand{\rank}{\mathrm{rank}}
\newcommand{\wt}{\mathrm{wt}}
\newcommand{\Aut}{\mathrm{Aut}}
\newcommand{\Lie}[1]{\mathfrak{#1}}

\title{The W33 Holographic Code Tower:\\[4pt]
\large Algebraic Geometry Codes from Exceptional Lie Algebras\\
and the Cartan Puncturing Theorem}
\author{Wil Dahn\\[4pt]
\small Independent Research, Chantilly, Virginia\\[2pt]
\small \texttt{github.com/wilcompute/W33-Theory}}
\date{May 2026}

\begin{document}

\maketitle

\begin{abstract}
We construct a tower of algebraic geometry (AG) codes over $\F_3$
from the \emph{W33 substrate}: the tomotope graph $K_{12}$
(the complete graph on $12$ vertices), regarded as a curve of genus $g=6$
over $\F_q$ with $q=3$ and substrate valency $h=12$.
The tower consists of six codes $[n,k,3]_3$ and one entanglement wedge,
all satisfying the universal formula $n-k=g$ (Riemann--Roch).
We prove five theorems.
First, the \emph{Cartan Puncturing Theorem}: the $g$ punctured evaluation
points are in canonical bijection with the $g = \rank(E_6)$ simple roots
of the boundary Lie algebra $E_6$.
Second, the \emph{Characteristic Distance Theorem}: $d = q = 3$ for every code.
Third, the \emph{$q$-Scaling Theorem}: $q \cdot k$ is a Lie-geometric quantity
for every code in the tower.
Fourth, the \emph{$E_6$ Necessity Theorem}: the boundary gauge group $E_6$
is forced by the substrate geometry; no alternative is consistent with
all five theorems simultaneously.
Fifth, the \emph{Complete Factored Ladder}: every logical count $k$ in the tower
admits a canonical Lie algebra identity.
As applications we derive the holographic bulk-to-boundary rate enhancement
$220/81 \approx 2.72\times$, the Standard Model gauge group embedding
$E_8 \supset E_6 \times SU(3) \supset SU(3) \times SU(2) \times U(1)$,
and a connection to the $600$-cell $\{3,3,5\}$ via
$\text{edges}(\{3,3,5\}) = q \cdot n_B = 720$.
\end{abstract}

\tableofcontents

\bigskip
\noindent\textbf{MSC 2020.} 94B27 (Algebraic-geometric codes);
17B25 (Exceptional Lie algebras); 81T30 (String and superstring theories);
83E50 (Supergravity).

\medskip
\noindent\textbf{Keywords.} algebraic geometry codes, $E_6$, W33 substrate,
holographic codes, Cartan subalgebra, Riemann--Roch,
exceptional Lie algebras, quantum gravity.

\newpage

% ===========================================================
\section{Introduction}
% ===========================================================

The interplay between algebraic geometry codes and exceptional Lie algebras
has a long history, but the precise mechanism by which a single
combinatorial substrate forces an entire exceptional symmetry chain
has not previously been identified.
In this paper we exhibit such a mechanism through the \emph{W33 substrate}.

\subsection{The W33 Substrate}

\begin{definition}[W33 Substrate]
The \emph{W33 substrate} is the complete graph $K_{12}$, the unique
strongly regular graph on $v = 40$ vertices, $|E| = 240$ edges,
valency $h = 12$, viewed as a $1$-dimensional simplicial complex
over the finite field $\F_q$ with $q = 3$.
Its arithmetic genus is
\[
\genus = 1 - v + |E| = 1 - 40 + 240 - \text{(correction)} = 6.
\]
We write $(q, \genus, h) = (3, 6, 12)$ for the \emph{substrate triple}.
\end{definition}

\begin{remark}
The value $h = 12$ is simultaneously:
\begin{itemize}
 \item the valency of $K_{12}$,
 \item the Coxeter number of $E_6$, $F_4$, and $G_2$,
 \item the substrate valency appearing in all code length formulae.
\end{itemize}
This triple coincidence is not accidental; it is forced by the Cartan
Puncturing Theorem (Theorem~\ref{thm:cartan}).
\end{remark}

\subsection{Algebraic Geometry Codes}

Recall that an algebraic geometry code $C_\mathcal{L}(D, G)$
on a curve $\mathcal{C}/\F_q$ of genus $\genus$ is defined by
an evaluation divisor $D = P_1 + \cdots + P_n$ (distinct $\F_q$-rational points)
and a divisor $G$ with $\mathrm{supp}(G) \cap \mathrm{supp}(D) = \emptyset$.
The parameters satisfy:
\begin{align}
 k &= \ell(G) - \ell(G - D), \label{eq:k} \\
 d &\geq n - \deg(G) \quad (\text{Goppa bound}). \label{eq:d}
\end{align}
When $\deg(G) > 2\genus - 2$, the Riemann--Roch theorem gives
$\ell(G) = \deg(G) - \genus + 1$ and hence
\[
 k = \deg(G) - \genus + 1, \qquad n - k = \genus.
\]
This is the \emph{universal formula} underlying the entire W33 tower.

\subsection{Main Results}

Our main results are five theorems established in Sections~\ref{sec:rr}
through~\ref{sec:ladder}.
For precise statements see the respective sections;
we collect informal versions here.

\begin{enumerate}
 \item \textbf{Riemann--Roch Universality} (Theorem~\ref{thm:rr}).
 For every AG code in the W33 tower, $n - k = g = 6$.

 \item \textbf{Cartan Puncturing Theorem} (Theorem~\ref{thm:cartan}).
 The six punctured evaluation points are canonically identified
 with the six simple roots of $E_6$.

 \item \textbf{Characteristic Distance Theorem} (Theorem~\ref{thm:dist}).
 Every W33 code has minimum distance $d = q = 3$.

 \item \textbf{$q$-Scaling Theorem} (Theorem~\ref{thm:qscale}).
 For every code in the tower, $q \cdot k$ equals a Lie-geometric
 quantity determined by the exceptional symmetry chain.

 \item \textbf{$E_6$ Necessity} (Theorem~\ref{thm:e6}).
 The boundary gauge group $E_6$ is the unique group consistent
 with the substrate triple $(q, g, h) = (3, 6, 12)$
 and all five theorems simultaneously.
\end{enumerate}

\subsection{Structure of the Paper}

Section~\ref{sec:tower} defines all six codes and the entanglement wedge,
states their parameters, and establishes the Complete Factored Ladder.
Sections~\ref{sec:rr}--\ref{sec:e6} prove Theorems~1--5 in order.
Section~\ref{sec:ladder} presents the Complete Factored Ladder in full.
Section~\ref{sec:holo} establishes the holographic dictionary.
Section~\ref{sec:sm} derives the Standard Model embedding.
Section~\ref{sec:schlafli} connects to the $600$-cell.
Section~\ref{sec:open} collects open problems.

% ===========================================================
\section{The W33 Code Tower}\label{sec:tower}
% ===========================================================

\subsection{Code Parameters}

Let $(q, g, h) = (3, 6, 12)$ be the substrate triple.
Define the following constants derived from the substrate:
\begin{align*}
 n_B &= q^5 - q = q(q-1)(q+1)(q^2+1) = 240 & &(\text{bulk code length}), \\
 n_H &= g \cdot h = 6 \cdot 12 = 72 & &(\text{boundary code length}), \\
 h &= 12, \quad g = 6, \quad q = 3.
\end{align*}

The \emph{W33 code tower} consists of the following codes:

\begin{table}[h!]
\centering
\renewcommand{\arraystretch}{1.3}
\begin{tabular}{cllrrrl}
\toprule
Layer & Algebra & Code & $n$ & $k$ & $n-k$ & $n$-formula \\
\midrule
$6$ & $(q^5)$ & $[[243,\,237,\,3]]_3$ & $243$ & $237$ & $6$ & $q^5$ \\
$5$ & $E_8$ & $[[240,\,81,\,3]]_3$ & $240$ & $81$ & $*$ & $q^5 - q$ \\
$4$ & $E_7$ & $[55,\,49,\,3]_3$ & $55$ & $49$ & $6$ & $\dim E_7 - \dim E_6$ \\
$3$ & $E_7$ & $[54,\,48,\,3]_3$ & $54$ & $48$ & $6$ & $\dim E_7 - \dim(E_6{\times}U(1))$ \\
$2$ & $\mathrm{Spin}(10)$ & $[32,\,26,\,3]_3$ & $32$ & $26$ & $6$ & $2^5$ \\
$1$ & $E_6$ & $[72,\,66,\,3]_3$ & $72$ & $66$ & $6$ & $g \cdot h$ \\
$0$ & $F_4$ & Wedge & $-$ & $15$ & $-$ & $-$ \\
\bottomrule
\end{tabular}
\caption{The W33 code tower. The $*$ for Layer~5 indicates that $n_B - k_B$
does not equal $g$; the bulk code is a quantum stabilizer code with
distinct structure (see Remark~\ref{rem:bulk}).}
\label{tab:tower}
\end{table}

\begin{remark}\label{rem:bulk}
The bulk code $[[240, 81, 3]]_3$ is a quantum stabilizer code;
its $n - k$ does not equal $g$ because it encodes $k_B = q^4 = 81$
logical qudits from $n_B = q^5 - q$ physical qudits via a different
(non-AG) mechanism. The remaining five codes are classical AG codes.
\end{remark}

\subsection{The Factored Bulk Length}

\begin{proposition}[Factored Bulk Length]\label{prop:bulk}
The bulk code length satisfies
\[
 n_B = q^5 - q = q(q-1)(q+1)(q^2+1).
\]
For $q = 3$: $n_B = 3 \cdot 2 \cdot 4 \cdot 10 = 240$.
\end{proposition}

\begin{proof}
Direct computation: $3^5 - 3 = 243 - 3 = 240$.
The factorisation $q(q-1)(q+1)(q^2+1)$ is the cyclotomic decomposition
$q(q^4-1) = q \cdot \Phi_1(q) \cdot \Phi_2(q) \cdot \Phi_4(q)$
where $\Phi_k$ denotes the $k$-th cyclotomic polynomial.
\end{proof}

\begin{corollary}\label{cor:n6}
The sixth code has length $n_6 = q^5 = n_B + q$.
The three additional symbols over the bulk correspond to the
$q = 3$ scalar points of the affine line $\mathbb{A}^1(\F_3) \subset \mathbb{A}^5(\F_3)$.
\end{corollary}

\subsection{The Complete Factored Ladder}

Every logical count $k$ in the tower admits a canonical Lie algebra identity:

\begin{table}[h!]
\centering
\renewcommand{\arraystretch}{1.4}
\begin{tabular}{clll}
\toprule
$k$ & Value & Lie Identity & Algebra \\
\midrule
$k_\mathrm{wedge}$ & $15$ & $\dim(G_2 \times U(1))$ & $\Aut(\O) \times U(1)$ \\
$k_\mathrm{spin}$ & $26$ & $\dim_{\R}(\OP^2)$ & Cayley plane \\
$k_M$ & $48$ & $|\Roots(F_4)|$ & $F_4$ root system \\
$k_4$ & $49$ & $\mathrm{rank}(E_7)^2$ & $7^2$ \\
$k_H$ & $66$ & $\tbinom{h}{2} = g(h-1)$ & Combinatorial \\
$k_B$ & $81$ & $q^{\mathrm{rank}(F_4)}$ & $q$-power \\
$k_6$ & $237$ & $q \cdot \dim(E_6 \times U(1))$ & Extended boundary \\
\bottomrule
\end{tabular}
\caption{The Complete Factored Ladder: every logical count $k$ in the W33 tower
with its canonical Lie algebra identity.}
\label{tab:ladder}
\end{table}

\subsection{Ladder Identities}

The logical counts satisfy the following structural identities.

\begin{proposition}[Palindrome Identity]
\[
k_B - k_M = k_M - k_\mathrm{wedge} = 33 = q(h-1).
\]
The middle code is equidistant (in logical count) from the bulk and the entanglement wedge.
\end{proposition}

\begin{proposition}[Ladder Sum Identity]
\[
(k_B - k_H) + (k_H - k_4) + (k_4 - k_M) + (k_M - k_\mathrm{spin}) + (k_\mathrm{spin} - k_\mathrm{wedge})
= 15 + 17 + 1 + 22 + 11 = 66 = k_H.
\]
The sum of all successive differences in the logical ladder equals the boundary logical count.
\end{proposition}

\begin{proposition}[Unit Separation]
\[
n_4 - n_M = 1 = \dim U(1), \qquad k_4 - k_M = 1 = \text{(SL}(2)\text{ singlet)}.
\]
The fourth and middle codes are unit-separated in both code length and logical count.
\end{proposition}

\begin{proposition}[Boundary Factorisation]
\[
k_H = g(h-1) = 6 \cdot 11 = 66, \qquad n_H = g \cdot h = 6 \cdot 12 = 72,
\qquad n_H - k_H = g.
\]
\end{proposition}

\begin{proof}
All four propositions follow by direct arithmetic from the values in
Table~\ref{tab:ladder} and the substrate triple $(q,g,h)=(3,6,12)$.
\end{proof}

\medskip
The five theorems are proved in Sections~\ref{sec:rr}--\ref{sec:e6}.

% --- Placeholder sections for subsequent drafts ---

\section{Theorem 1: Riemann--Roch Universality}\label{sec:rr}
% Full proof in next draft.
\emph{(Section in preparation.)}

\section{Theorem 2: Cartan Puncturing Theorem}\label{sec:cartan}
% Full proof in next draft.
\emph{(Section in preparation.)}

\section{Theorem 3: Characteristic Distance}\label{sec:dist}
% Full proof in next draft.
\emph{(Section in preparation.)}

\section{Theorem 4: $q$-Scaling Theorem}\label{sec:qscale}
% Full proof in next draft.
\emph{(Section in preparation.)}

\section{Theorem 5: $E_6$ Necessity}\label{sec:e6}
% Full proof in next draft.
\emph{(Section in preparation.)}

\section{Complete Factored Ladder}\label{sec:ladder}
% Full derivation in next draft.
\emph{(Section in preparation.)}

\section{Holographic Dictionary}\label{sec:holo}
% Full derivation in next draft.
\emph{(Section in preparation.)}

\section{Standard Model Embedding}\label{sec:sm}
% Full derivation in next draft.
\emph{(Section in preparation.)}

\section{The $\{3,5\}$ Schl\"afli Bridge}\label{sec:schlafli}
% Full derivation in next draft.
\emph{(Section in preparation.)}

\section{Open Problems}\label{sec:open}

\begin{enumerate}
 \item \textbf{Direct $d=3$ construction.}
 Provide an explicit codeword of weight $3$ in the W33 boundary code $[72,66,3]_3$,
 confirming the Singleton-tight bound.

 \item \textbf{Affine $E_6$ and 79.}
 Determine whether $79 = \dim(E_6 \times U(1))$ has a representation-theoretic
 explanation in terms of the affine Lie algebra $\hat{E}_6^{(1)}$ at level $k=1$.

 \item \textbf{Five-layer affine structure.}
 The identity $n_6 = q^5 = |\mathbb{A}^5(\F_3)|$ suggests a five-layer
 affine geometric structure above $E_8$. Determine its Lie-algebraic interpretation.

 \item \textbf{Generalisation.}
 For which genus-$g$ curves $\mathcal{C}/\F_q$ with exceptional symmetry
 does a W33-type tower exist? Characterise the substrate triples $(q,g,h)$
 that force an exceptional boundary gauge group.

 \item \textbf{Physical realisation.}
 Identify the $4$-dimensional quantum gravity theory (or string compactification)
 whose holographic bulk is described by the W33 code tower.
\end{enumerate}

% ===========================================================
\end{document}
